\section{Conclusão}

Este trabalho apresentou o Wumpus Verse, um protótipo de plataforma web desenvolvido para apoiar a criação, execução e avaliação de agentes inteligentes no ambiente do Mundo do Wumpus. A proposta surgiu da necessidade de disponibilizar uma ferramenta acessível e integrada que permitisse a experimentação com agentes inteligentes sem exigir dos usuários conhecimentos avançados de programação ou configuração de ambientes de teste.

A plataforma foi projetada para oferecer recursos que abrangem todo o ciclo de experimentação, incluindo a criação de mundos personalizados, a configuração de diferentes tipos de agentes, a execução de simulações, o armazenamento de resultados e a geração de estatísticas de desempenho. Além disso, sua implementação em ambiente web amplia a acessibilidade da ferramenta, permitindo sua utilização para a realização de experimentação em lote e padronização de métricas.

Os testes realizados demonstraram o funcionamento adequado das principais funcionalidades do sistema, evidenciando a integração entre os módulos desenvolvidos e a capacidade da plataforma de fornecer um ambiente consistente para a realização de experimentos envolvendo agentes inteligentes.

Como perspectivas futuras, pretende-se ampliar as funcionalidades da plataforma de modo a torná-la ainda mais flexível e personalizável. Uma das principais evoluções planejadas está relacionada ao modelo de construção dos ambientes. Atualmente, embora os mundos possam assumir diferentes formatos por meio da adição de salas em posições arbitrárias, sua estrutura ainda permanece limitada a uma malha bidimensional. Como extensão, pretende-se permitir a criação de conexões manuais entre salas, de forma semelhante a diagramas ou grafos, possibilitando a construção de ambientes com topologias mais complexas e menos restritas ao modelo tradicional do Mundo de Wumpus.

Também estão previstas melhorias nos mecanismos de execução das simulações, oferecendo maior controle sobre parâmetros experimentais e ampliando as opções de configuração disponíveis aos usuários. Dessa forma, espera-se proporcionar maior flexibilidade para a realização de diferentes tipos de experimentos e análises.

Outra funcionalidade de interesse é o suporte a múltiplos agentes atuando simultaneamente em um mesmo ambiente. Essa extensão permitirá a investigação de cenários cooperativos e competitivos, ampliando significativamente as possibilidades de estudo e avaliação de estratégias de tomada de decisão.

Por fim, prevê-se a incorporação de recursos voltados à interação entre usuários, como compartilhamento de mundos e agentes, sistemas de classificação e mecanismos de colaboração.

Com essas evoluções, espera-se consolidar o Wumpus Verse como uma plataforma cada vez mais completa e acessível para atividades de ensino, pesquisa e experimentação envolvendo agentes inteligentes.